\textbf{Data.} Each graph is a random spanning tree (each node attached to a uniformly random earlier node in a random order) plus independent Erd\H{o}s--R\'enyi edges, undirected, weights $U(0.2,1)$, uniformly random source. Training: $n\sim U\{4,\dots,8\}$ per batch, ER probability $p=0.4$, fresh graphs every step (no fixed training set). Test: 200 graphs per cell, $n\in\{8,16,32,64\}$ and two families: \emph{sparse} ($p=3/(n-1)$, constant expected extra degree) and \emph{dense} ($p=0.4$, degree grows with $n$). Cells are written S$n$ and D$n$. The test graphs of seed $k$ are drawn from a random stream seeded with $10000+k$, which is separate from the training stream (seed $k$) and does not depend on the system. Each seed therefore has its own test set, and at a given seed every system, including the ablation and sweep runs, is evaluated on the same 200 graphs per cell; the paired tests pair on this seed.

\textbf{Metrics.} Primary: mean absolute error (MAE) of the predicted final distance over all nodes, cell S64, averaged over seeds. Also relative error $\sum|\hat d-d|/\sum d$ (rS64, rD64; predicting zero everywhere gives exactly 1), and the fraction of nodes within 0.05. Because the distance scale differs between families (dense graphs have short paths), absolute errors are not comparable across families; the relative error is. A seed is counted as a \emph{failure} in a cell if its MAE exceeds 1.0 or is non-finite (``div.'' in the tables means a non-finite value, i.e.\ the network output overflowed). The threshold, the medians and the failure counts were introduced after seeing the first results; they are descriptive and not part of the registered tests.

\textbf{Protocol.} Five seeds (0--4) for all five main systems; three seeds (0--2) for ablations. Adam, learning rate $3\times10^{-3}$ with cosine decay, 1000 steps, batch 32 graphs, gradient-norm clip 1, identical for every system. We did \emph{no} hyperparameter tuning (the defaults were fixed before any run) and no test-size data was used for any choice; a validation MAE on the training distribution is logged but unused. All systems use the same data generators, metric code and number of message-passing steps ($T=n-1$ at train and test, except in the step sweep of Section~\ref{sec:ablations}). Hardware: shared laptop CPU, 2 threads; the slowest run took about two minutes. The code, run log and per-seed metrics are in the repository.

\textbf{Hypotheses and registered tests} (\texttt{proposal.md}, written before the runs). \textbf{H1} max+steps has lower S64 MAE than sum+steps; \textbf{H2} step supervision lowers the max-MPNN's S64 MAE; \textbf{H3} the MLP is worse than every MPNN for $n\ge16$ (tested at S16 and S64); \textbf{H4} under growing degree the per-seed ratio D64/D8 of sum+steps exceeds that of max+steps. H1--H3 are decided by the Welch and paired tests of \texttt{rh compare} on seed-level MAE against the reference system max+steps, and count as not supported if the sign is wrong or the difference is not significant at $0.05$. No test was registered for H4. Everything else we report is post hoc and marked as such: the S16 tests for H1 and H2, the final-only pair for H4, and the $\times4$ column of the step sweep (the registered sweep was $\{0.5,1,2\}$).
